% =============================================================================
%  ISM COMPANION READER  —  SEML Session 6
%  EDA, Model Registry & Serving Patterns
%  Built from templates/companion.tex (gold-standard faithful)
% =============================================================================

\documentclass[a4paper,11pt]{article}

% ---- Encoding & fonts -------------------------------------------------------
\usepackage[T1]{fontenc}
\usepackage[utf8]{inputenc}
\usepackage{lmodern}
\usepackage{microtype}

% ---- Geometry ---------------------------------------------------------------
\usepackage[a4paper,margin=1in,headheight=15pt,headsep=14pt]{geometry}

% ---- Math -------------------------------------------------------------------
\usepackage{amsmath,amssymb}

% ---- Graphics, tables, colour ----------------------------------------------
\usepackage{graphicx}
\usepackage{booktabs}
\usepackage[table,svgnames]{xcolor}

% ---- Callout boxes ----------------------------------------------------------
\usepackage[most]{tcolorbox}
\tcbuselibrary{skins,breakable}

% ---- Tikz (specimen fallback) -----------------------------------------------
\usepackage{tikz}

% ---- Glyphs -----------------------------------------------------------------
\usepackage{pifont}
% Use pifont glyphs unconditionally (bbding command names vary across installs).
\newcommand{\glyphHand}{\ding{43}}%
\newcommand{\glyphPencil}{\ding{46}}%

% ---- Lists & spacing --------------------------------------------------------
\usepackage{enumitem}
\setlist{leftmargin=1.5em,topsep=4pt,itemsep=3pt}

% Steps list: keeps "Step N." INSIDE the callout box -------------------------
\newlist{steps}{enumerate}{1}
\setlist[steps]{label=\textbf{Step \arabic*.},
  leftmargin=4.4em, labelwidth=3.8em, labelsep=0.6em, labelindent=0pt,
  align=left, topsep=5pt, itemsep=6pt}
\usepackage{parskip}

% ---- Headers / footers ------------------------------------------------------
\usepackage{fancyhdr}

% ---- Section styling --------------------------------------------------------
\usepackage{titlesec}

% ---- Links + URL line-breaking ----------------------------------------------
\usepackage[hidelinks]{hyperref}
\usepackage{xurl}
\hypersetup{breaklinks=true}

% ---- adjustbox for wide tables ----------------------------------------------
\usepackage{adjustbox}

% =============================================================================
%  COLOURS
% =============================================================================
\definecolor{navy}{HTML}{21355E}
\definecolor{navyrule}{HTML}{21355E}
\definecolor{inktext}{HTML}{1A202C}
\definecolor{mutedtext}{HTML}{6B7280}

\definecolor{intuFrame}{HTML}{2C5AA0}   \definecolor{intuFill}{HTML}{F4F7FC}
\definecolor{everyFrame}{HTML}{8A5A1E}  \definecolor{everyFill}{HTML}{FBF1E3}
\definecolor{workFrame}{HTML}{2E7D52}   \definecolor{workFill}{HTML}{F1F8F3}
\definecolor{watchFrame}{HTML}{B23A48}  \definecolor{watchFill}{HTML}{FBF1F2}
\definecolor{keyFrame}{HTML}{6A4C93}    \definecolor{keyFill}{HTML}{F4F0F9}

\color{inktext}
\hypersetup{linkcolor=navy,urlcolor=navy,citecolor=navy}

% =============================================================================
%  RUNNING HEADER / FOOTER
% =============================================================================
\newcommand{\CourseShort}{SEML Companion}
\newcommand{\SessionNo}{6}
\newcommand{\LectureTitle}{EDA, Model Registry \& Serving Patterns}

\pagestyle{fancy}
\fancyhf{}
\fancyhead[L]{\footnotesize\color{mutedtext}\textsf{\CourseShort}}
\fancyhead[R]{\footnotesize\color{mutedtext}\textsf{Session \SessionNo\ \textperiodcentered\ \LectureTitle}}
\fancyfoot[C]{\footnotesize\color{mutedtext}\thepage}
\renewcommand{\headrulewidth}{0.4pt}
\renewcommand{\footrulewidth}{0pt}
\renewcommand{\headrule}{\hbox to\headwidth{\color{navyrule!55}\leaders\hrule height \headrulewidth\hfill}}

% =============================================================================
%  SECTION HEADINGS
% =============================================================================
\titleformat{\section}
  {\sffamily\Large\bfseries\color{navy}}
  {\thesection}{0.8em}{}
  [{\vspace{2pt}\color{navy!70}\titlerule[0.8pt]}]
\titleformat{\subsection}
  {\sffamily\large\bfseries\color{navy!92}}
  {\thesubsection}{0.7em}{}
\titlespacing*{\section}{0pt}{16pt}{8pt}
\titlespacing*{\subsection}{0pt}{12pt}{4pt}

% =============================================================================
%  PILL-TAB CALLOUTS
% =============================================================================
\tcbset{
  housecallout/.style 2 args={
    enhanced, breakable,
    colback=#2, colframe=#1,
    boxrule=0.6pt, arc=2.4pt, outer arc=2.4pt,
    left=10pt, right=10pt, top=9pt, bottom=9pt,
    before skip=11pt, after skip=11pt,
    fontupper=\normalsize,
    coltitle=white,
    fonttitle=\bfseries\sffamily\small,
    attach boxed title to top left={xshift=6mm,yshift=-3mm},
    boxed title style={
      colback=#1, colframe=#1,
      rounded corners, arc=3pt,
      boxrule=0pt, left=6pt, right=6pt, top=2pt, bottom=2pt
    }
  }
}

\newtcolorbox{intuition}{%
  housecallout={intuFrame}{intuFill},
  title={\raisebox{-0.05ex}{\glyphHand}\,\ The intuition}}

\newtcolorbox{everyday}{%
  housecallout={everyFrame}{everyFill},
  title={\raisebox{-0.05ex}{$\bigstar$}\,\ Everyday picture}}

\newtcolorbox{worked}[1]{%
  housecallout={workFrame}{workFill},
  title={\raisebox{-0.05ex}{\glyphPencil}\,\ Worked example: #1}}

\newtcolorbox{watchout}{%
  housecallout={watchFrame}{watchFill},
  title={\raisebox{-0.05ex}{\ding{55}}\,\ Watch out}}

\newtcolorbox{keytake}{%
  housecallout={keyFrame}{keyFill},
  title={\raisebox{-0.05ex}{\ding{51}}\,\ Key takeaway}}

% =============================================================================
%  TITLE BANNER
% =============================================================================
\providecommand{\addfontfeatureslike}[1]{#1}
\newcommand{\titlebanner}[4]{%
  \begin{tcolorbox}[
    enhanced, breakable=false,
    colback=navy, colframe=navy,
    boxrule=0pt, arc=3pt, outer arc=3pt,
    left=16pt, right=16pt, top=16pt, bottom=16pt,
    before skip=2pt, after skip=14pt,
    width=\linewidth ]
    {\sffamily\footnotesize\bfseries\color{white!85}%
      \MakeUppercase{\addfontfeatureslike{#1}}}\par
    \vspace{6pt}
    {\sffamily\bfseries\fontsize{30}{34}\selectfont\color{white} #2\par}
    \vspace{5pt}
    {\sffamily\large\color{white!88} #3\par}
    \vspace{9pt}
    {\color{white!55}\rule{\linewidth}{0.4pt}}\par
    \vspace{6pt}
    {\sffamily\footnotesize\color{white!82} #4\par}
  \end{tcolorbox}%
}

% =============================================================================
%  \housefig
% =============================================================================
\newcommand{\housefig}[2]{%
  \begin{figure}[htbp]
    \centering
    \includegraphics[width=\linewidth]{#1}%
    \caption{#2}
  \end{figure}%
}

\newcommand{\vocab}[1]{\textbf{\color{navy}#1}}

% =============================================================================
\begin{document}
% =============================================================================

\thispagestyle{fancy}

\titlebanner
  {Software Engineering for Machine Learning \textperiodcentered\ Companion Reader}
  {EDA, Model Registry \& Serving Patterns}
  {Architectural patterns that keep ML systems reliable and scalable}
  {Session 6 \textperiodcentered\ Read alongside the lecture slides \textperiodcentered\ Every slide example worked out in full}

% ---- How to use this companion ----------------------------------------------
\noindent\textbf{How to use this companion.}
The slides give you the formulas; this reader gives you the \emph{why}.
Every slide concept is expanded into a full, plain-English explanation.
Every slide example is worked out step by step.
Colour-coded boxes guide you through each concept.
The \textcolor{intuFrame}{\textbf{blue}} ``intuition'' box states the core idea
in one breath.
The \textcolor{everyFrame}{\textbf{amber}} ``everyday picture'' box ties it to
real life before any math.
The \textcolor{workFrame}{\textbf{green}} ``worked example'' box solves a
problem with real numbers and no skipped steps.
The \textcolor{watchFrame}{\textbf{red}} ``watch out'' box flags the common
mistake.
The \textcolor{keyFrame}{\textbf{purple}} ``key takeaway'' box is the one
sentence to remember. Read in order alongside the slides.

\medskip

% =============================================================================
\section{What this session covers}
% =============================================================================

Session~6 asks one practical question: \emph{how do we build ML systems that
are reliable, scalable, and easy to maintain?}
The answer comes in three architectural patterns.

\begin{enumerate}
  \item \textbf{Event-Driven Architecture (EDA)} --- how services talk to each
        other without being tightly joined. Components publish events; other
        components listen and react.
  \item \textbf{Model Registry Pattern} --- a central store for every version
        of every ML model. It bridges experimentation and production.
  \item \textbf{Batch vs.\ Real-Time Serving} --- two ways to deliver
        predictions. Batch runs offline on a schedule; real-time answers in
        milliseconds.
\end{enumerate}

Each pattern solves a recurring problem. By the end you should be able to say
\emph{why} you would pick one over the other in a real project.

% =============================================================================
\section{Event-Driven Architecture: services that talk without being glued}
% =============================================================================

Picture a big software system with many moving parts: a web server,
a recommendation engine, a notification service, a logging pipeline.
If every part calls every other part directly, the system becomes a tangled
web. Change one part and everything else may break.
\vocab{Event-Driven Architecture (EDA)} is the pattern that cuts those direct
connections. Instead of calling each other directly, services produce and
consume \vocab{events} --- small messages that say ``something just happened.''

\begin{intuition}
In EDA, services do not talk to each other directly. They shout into a shared
room. Any service that cares can listen. Services that do not care can ignore
the shout entirely. The room is the \textbf{event bus} or
\textbf{message broker}.
\end{intuition}

\begin{everyday}
Think of a newspaper publisher. The publisher drops copies at a distribution
centre. Subscribers pick up only the sections they care about: sport,
business, or technology. The publisher has no idea who is reading.
Readers have no idea who else is reading. They are \emph{decoupled}.

In EDA, the publisher is a \textbf{producer}, the distribution centre is the
\textbf{event bus}, and the readers are \textbf{consumers} (subscribers).
\end{everyday}

\subsection{The two roles: producers and consumers}

A \vocab{producer} (also called a \emph{publisher}) is any service that
generates an event. It might be a user clicking a button, an external API
sending data, or an internal timer firing at midnight. The producer sends the
event to the event bus and then forgets about it.

A \vocab{consumer} (also called a \emph{subscriber}) is any service that
registered its interest in a type of event. When the bus receives a matching
event, it notifies the consumer. Many consumers can listen to the same event.

\subsection{The event bus and message brokers}

The \vocab{event bus} is the middleman. Real-world tools that play this role:

\begin{itemize}
  \item \textbf{Apache Kafka} --- high-throughput, fault-tolerant, built for
        streaming data at scale.
  \item \textbf{RabbitMQ} --- flexible routing, good for task queues.
  \item \textbf{Amazon SQS / Azure Service Bus} --- managed cloud services;
        no infrastructure to run yourself.
\end{itemize}

These tools store events even when a consumer is temporarily offline.
When the consumer comes back, it catches up from where it left off.

\subsection{Why EDA matters: the four problems it solves}

The lecture lists four reasons to use EDA.

\begin{enumerate}
  \item \textbf{Async communication.} Events happen at different times.
        EDA lets services process events when they are ready, not in lockstep.
  \item \textbf{Decoupling.} Producers and consumers are independent.
        You can add a new consumer without touching the producer at all.
  \item \textbf{Horizontal scalability.} Need to handle more events?
        Add more consumer instances. The bus handles the distribution.
  \item \textbf{Reliability.} The broker stores events persistently.
        A consumer crash does not lose data; it resumes later.
\end{enumerate}

\begin{worked}{an e-commerce ML pipeline using EDA}
An online store uses EDA to feed its recommendation model.

\begin{steps}
  \item \textbf{User action triggers an event.}
        A user clicks ``Add to cart.'' The web server (producer) publishes
        one event to the bus:
        \texttt{\{user: 42, product: "headphones", event: "add\_to\_cart"\}}.
  \item \textbf{The bus stores and routes the event.}
        Kafka writes the event to a topic called
        \texttt{user-interactions}. Any subscriber to that topic gets a copy.
  \item \textbf{Three consumers react in parallel.}
        Consumer~1 (ML pipeline) updates the user's feature vector.
        Consumer~2 (analytics) increments a counter for headphones sold today.
        Consumer~3 (inventory) checks stock levels.
        None of the three knows about the others.
  \item \textbf{No consumer changes the producer.}
        The web server code never changed. Three consumers were added
        without touching the producer. This is the decoupling payoff.
  \item \textbf{Reliability check.}
        If Consumer~1 crashes during step~3, Kafka keeps the event.
        When Consumer~1 restarts, it reads from its saved offset and
        processes the event. No data is lost.
\end{steps}
The ML pipeline, analytics, and inventory all stay in sync from one event,
with no direct calls between any of them.
\end{worked}

\housefig{figures/fig_eda_pubsub.png}{%
  The publish-subscribe flow in EDA. Producers push events to the broker.
  The broker routes each event to all registered consumers. No producer
  knows which consumers exist --- that is the decoupling.}

\begin{watchout}
\begin{itemize}
  \item \textbf{Eventual consistency.} Consumers react async, so they may
        lag behind the producer. Do not assume all consumers have the latest
        state at any given instant.
  \item \textbf{Ordering is not free.} Events may arrive out of order.
        Kafka guarantees order within a partition, but you must set up
        partition keys carefully.
  \item \textbf{Debugging is harder.} A bug in a consumer is invisible to
        the producer. You need distributed tracing to follow an event's path.
  \item \textbf{Overhead for simple systems.} Two services talking? A direct
        API call is simpler. EDA adds complexity; use it only when you need it.
\end{itemize}
\end{watchout}

\noindent\textbf{Where this shows up in ML.}
An ML system is a \emph{pipeline}: data collection, feature engineering,
training, serving, monitoring. EDA lets each stage be an independent service.
A new batch of training data triggers a \texttt{data-available} event.
The feature store consumes it and fires \texttt{features-ready}.
The training job starts. When done, a \texttt{model-trained} event fires.
The serving layer picks it up and swaps the model. Each stage is replaceable.
Tools like \textbf{Apache Kafka} and \textbf{MLflow} are often combined
in exactly this pattern.

\begin{keytake}
EDA decouples services through an event bus. Producers publish; consumers
subscribe. No direct connections means each part can change, scale, and fail
independently --- the key property for a reliable ML system.
\end{keytake}

% =============================================================================
\section{Design Patterns for ML Systems: why patterns exist}
% =============================================================================

A \vocab{design pattern} is a named, reusable solution to a problem that comes
up again and again. Patterns are not code you copy; they are templates you
adapt. In ML engineering, patterns solve recurring challenges in three areas:
data preparation, model training, and production serving.

\begin{intuition}
A design pattern is a recipe. A recipe does not hand you a finished cake.
It gives you a reliable sequence of steps for a class of cakes.
You adapt it to your oven and your ingredients. ML design patterns do the
same: they give you a proven structure for a recurring engineering problem.
\end{intuition}

Two patterns appear in almost every production ML project.
They are the \textbf{Model Registry Pattern} and the
\textbf{Batch vs.\ Real-Time Serving Pattern}.

% =============================================================================
\section{The Registry Pattern: a controlled lookup table}
% =============================================================================

Before the ML-specific model registry, it helps to see the general
\vocab{registry pattern} from software engineering.

A registry is a ``central hub to dynamically register, locate, and retrieve
class instances or services.'' Think of a phone book. You do not need to know
someone's address; you look up their name. The book is the registry.

\begin{everyday}
Imagine a hotel concierge desk. Every guest service (restaurant, taxi, spa)
is listed there. If a guest wants a taxi, they ask the concierge, not the
street. If a new restaurant opens, it registers with the concierge.
The concierge is the registry: a single, up-to-date directory that both
parties trust. Services find each other through the centre, not directly.
\end{everyday}

In ML, the registry pattern applies to \emph{machine learning artifacts}:
trained models, datasets, and experiment runs. The \vocab{ML Model Registry}
is ``a centralised repository that tracks, versions, manages, and governs ML
models throughout their lifecycle.''

\begin{keytake}
The registry pattern provides one authoritative place to register something
once. Others can find it by name, version, or tag --- without knowing where
it physically lives.
\end{keytake}

% =============================================================================
\section{Version Control: from source code to ML artifacts}
% =============================================================================

You already know that \vocab{Git} tracks changes to source code.
The model registry does the same job, but for ML artifacts.
Understanding the gap between them explains why a separate tool is needed.

\begin{intuition}
Git stores files: code, config, and documentation.
An ML registry stores \emph{trained models}: large binary files.
It also stores the metadata that makes them reproducible ---
hyperparameters, evaluation metrics, the dataset used, and the run ID.
Storing a 2~GB weights file in Git is slow and wasteful.
A registry is built for that job.
\end{intuition}

\subsection{Three generations of version control}

Version control tools evolved in three generations. The table below shows
the key jump at each step.

\medskip
\noindent
\begin{adjustbox}{max width=\linewidth}
\begin{tabular}{@{}llll@{}}
\toprule
\textbf{Generation} & \textbf{Network} & \textbf{Operations} & \textbf{Example tools} \\
\midrule
First (software)  & None          & One file at a time  & RCS, SCCS \\
Second (software) & Centralised   & Multi-file          & CVS, Subversion \\
Third (software)  & Distributed   & Changesets          & Git \\
Third (ML)        & Distributed + artifact tracking & Models, data, experiments & MLflow, DVC, W\&B \\
\bottomrule
\end{tabular}
\end{adjustbox}
\medskip

The ML third generation adds \emph{artifact tracking} on top of Git's
distributed model. It stores code \emph{and} the large binary outputs of
training (model weights). It also stores metadata linking each output to
its inputs: data version, hyperparameters, and metrics.

\housefig{figures/fig_vcs_generations.png}{%
  The four generations of version control, one per column.
  The ML generation adds artifact tracking on top of Git's
  distributed approach --- that is the new capability.}

\begin{worked}{what Git stores vs.\ what a model registry stores}
A team builds a sentiment classifier. They use both Git and MLflow.

\begin{steps}
  \item \textbf{What goes into Git.}
        The training script \texttt{train.py}, the requirements file
        \texttt{requirements.txt}, and \texttt{config.yaml} go into Git.
        These are text files. The team can see every code change over time.
  \item \textbf{What goes into the ML registry.}
        After training run~42, MLflow stores several things.
        It stores the trained model weights (450~MB).
        It stores the hyperparameters:
        \texttt{lr=0.001}, \texttt{epochs=10}, \texttt{batch=32}.
        It stores the validation accuracy (0.912) and the Git commit hash
        of \texttt{train.py} so the run is reproducible.
  \item \textbf{The link between the two systems.}
        MLflow records the Git commit hash alongside the model weights.
        If someone asks ``which code produced this model?'' they look up
        the hash and check out that commit in Git. The two systems are
        linked, not merged.
  \item \textbf{Why not just Git?}
        Storing a 450~MB weights file in Git makes every clone slow and
        every push expensive. Git also does not know what
        ``validation accuracy'' means; MLflow does.
\end{steps}
The two tools are complementary: Git for code, registry for artifacts.
\end{worked}

% =============================================================================
\section{Model Registry Pattern in MLOps}
% =============================================================================

In MLOps, the model registry is ``the definitive source of truth bridging
the gap between model experimentation and production serving.''
The registry is the \emph{handover point}.
It sits between the data scientist (who experiments) and the engineer
(who deploys). That hand-off is its most important job.

\begin{intuition}
Without a registry, deploying a model looks messy.
A data scientist emails a weights file to an engineer.
The engineer copies it to a server, hoping it is the right version.
With a registry, the data scientist registers the model.
The engineer promotes it to the \emph{production} alias. The serving layer
picks it up automatically. The whole process is auditable.
\end{intuition}

\subsection{What a model registry stores}

A model registry stores a set of linked artifacts for each model version.

\begin{itemize}
  \item \textbf{Trained model weights} --- the binary file(s) from training.
  \item \textbf{Datasets or dataset references} --- which data was used, or
        a pointer to it.
  \item \textbf{Hyperparameters} --- every setting used during training.
  \item \textbf{Metrics} --- accuracy, F1, loss, or any evaluation number.
  \item \textbf{Experiment logs} --- the training run that produced the model.
  \item \textbf{Lineage information} --- which code, data, and run produced
        this exact artifact. This is what makes a model \emph{reproducible}.
\end{itemize}

\subsection{MLflow Model Registry: the four key features}

\vocab{MLflow} is the most widely used open-source tool for ML experiment
tracking and model registry. Its registry has four main features.

\begin{enumerate}
  \item \textbf{Version control.} Every registration creates a new version
        number. You can compare versions, roll back, or run two versions
        in parallel for A/B testing.
  \item \textbf{Model lineage and traceability.} Each version links back to
        the MLflow run that produced it. From the run you can find the code,
        data, and hyperparameters.
  \item \textbf{Production-ready workflows.} MLflow supports \emph{aliases}
        (e.g.\ \texttt{production}, \texttt{staging}, \texttt{champion}) and
        \emph{tags} (e.g.\ \texttt{approved=true}).
        To promote a model, set its alias.
        The serving layer reads the alias, not a version number.
        It picks up the new model automatically.
  \item \textbf{Governance and compliance.} The registry keeps a full audit
        log of every promotion, annotation, and access event. This matters in
        regulated industries where you must prove which model made which
        decision at what time.
\end{enumerate}

\housefig{figures/fig_model_registry.png}{%
  The model lifecycle in an ML registry. A model moves left to right as it
  gains confidence. The registry tracks every transition with a timestamp
  and a user ID.}

\begin{worked}{promoting a model from staging to production in MLflow}
A team has trained a new churn-prediction model (version~3). They want to
replace the current production model (version~2).

\begin{steps}
  \item \textbf{Register the new model.}
        After training run~107, the data scientist calls:
        \texttt{mlflow.register\_model(run\_uri, "churn-model")}.
        MLflow creates \texttt{churn-model} version~3 with status
        \emph{None} (not yet promoted).
  \item \textbf{Validate in staging.}
        The team sets the alias:
        \texttt{client.set\_registered\_model\_alias(}
        \texttt{"churn-model", "staging", 3)}.
        The staging server reads the \texttt{staging} alias and loads
        version~3. QA tests pass: F1 = 0.87 vs.\ version~2's 0.82.
  \item \textbf{Promote to production.}
        With approval, the team updates the alias:
        \texttt{client.set\_registered\_model\_alias(}
        \texttt{"churn-model", "production", 3)}.
        The production endpoint reads the \texttt{production} alias. It
        now serves version~3. \textbf{No code change} was needed.
  \item \textbf{Keep version~2 as a fallback.}
        Version~2 stays in the registry with no alias. If version~3
        behaves badly, the team reassigns \texttt{production} back to
        version~2 in seconds.
  \item \textbf{Audit trail.}
        Every alias change is logged: who did it, when, and the previous
        value. This satisfies compliance requirements.
\end{steps}
The whole promotion took five steps and zero code deployments.
\end{worked}

\begin{watchout}
\begin{itemize}
  \item \textbf{The registry is not a serving layer.} It stores models; it
        does not run them. You still need a separate serving infrastructure
        that \emph{reads from} the registry.
  \item \textbf{Registry metrics are training metrics.}
        They do not tell you how the model behaves in production. Add a
        monitoring layer that tracks live performance separately.
  \item \textbf{Version numbers can mislead.}
        Version~5 is not always better than version~4. Always compare
        metrics; never assume a higher version number is better.
  \item \textbf{Lineage is only as good as your logging.}
        If a training run did not log its dataset version, the registry has
        an incomplete lineage record. Build logging into the training code
        from the start.
\end{itemize}
\end{watchout}

\noindent\textbf{Where this shows up in ML.}
Every serious MLOps platform uses a model registry at its core.
\textbf{MLflow Registry}, \textbf{Weights \& Biases Model Registry},
\textbf{DVC}, and cloud-native registries
(\textbf{AWS SageMaker}, \textbf{GCP Vertex AI}, \textbf{Azure ML Registry})
all implement this pattern. When you read about ``CI/CD for ML'' or
``model promotion workflows,'' the model registry is the component
that makes those workflows possible. Without it, there is no reliable way
to know which model is in production or to roll back when something goes wrong.

\begin{keytake}
The model registry is the single source of truth for ML models in production.
It tracks every version, every lineage link, and every promotion decision.
The serving layer reads from the registry by alias --- so promoting a model
requires no code change, just a registry update.
\end{keytake}

% =============================================================================
\section{Batch vs.\ Real-Time Serving: two ways to deliver predictions}
% =============================================================================

Once you have a model in the registry, you need to serve it: run it on new
data and return predictions. There are two fundamental patterns. The right
choice depends on how quickly the business needs the answer.

\begin{intuition}
Batch serving is like a bakery: bake a hundred loaves overnight, stock the
shelves in the morning, sell them through the day. Real-time serving is like
a food truck: each customer orders, you cook it fresh, and you hand it over
in two minutes. Both are valid. The question is what your customer needs.
\end{intuition}

\subsection{Batch prediction: processing data in bulk}

\vocab{Batch prediction} (also called \emph{offline prediction}) runs the
model on a large set of observations at once. It is a scheduled job, not a
live service. The results are stored in a database or file and consumed later.

The lecture gives a clear example: send a daily promotional email to customers
likely to buy a product. You do not need that list right now. A nightly job
can score all customers, save the top~1,000, and let the marketing team use
the list next morning.

\textbf{When to use batch prediction:}
\begin{itemize}
  \item \textbf{Customer segmentation} --- group customers by behaviour over
        the past week. No one needs the segments in milliseconds.
  \item \textbf{Product recommendations in emails} --- generate personalised
        recommendations for a daily newsletter. The email goes out once a day.
  \item \textbf{Risk reporting} --- compute credit risk scores for all loan
        applicants at end of business. The decision is made the next morning.
\end{itemize}

\begin{everyday}
A payroll system runs once a week. Every salary is computed in one batch job
on Friday night. No one needs their exact pay-to-date figure at 3~pm on
Wednesday. The batch job is fast, cheap, and perfectly correct for this use.
\end{everyday}

\subsection{Real-time (online) prediction: serving on demand}

\vocab{Online prediction} (also called \emph{real-time} or \emph{on-demand
inference}) deploys the model as a persistent web API. The service is always
running. When a request arrives, the model scores it and returns a prediction
in milliseconds.

The lecture's example is a credit card fraud check. The merchant needs
``approve'' or ``deny'' \emph{before the customer walks away}. A batch job
running tonight is useless. The model must be a live API that answers every
transaction in under 200~ms.

\textbf{When to use real-time prediction:}
\begin{itemize}
  \item \textbf{Fraud detection} --- score a transaction before it clears.
  \item \textbf{Dynamic pricing} --- adjust ride-share prices based on current
        supply and demand, updated every few seconds.
  \item \textbf{Spam filtering} --- decide if an incoming email is spam
        \emph{before} it reaches the user's inbox.
\end{itemize}

\begin{worked}{choosing the right serving pattern for three ML use cases}
For each use case, decide: batch or real-time?

\begin{steps}
  \item \textbf{Weekly newsletter recommendations.}
        The email goes out every Monday at 9~am. There is all weekend to
        score customers. Millions of customers need scoring.
        \textbf{Verdict: batch.}
        Run a Spark job on Sunday night. Store the top-5 recommendations
        per customer in a database. The email service reads the database
        on Monday. No live API needed.

  \item \textbf{Autocomplete for a search bar.}
        The user types one character. Suggestions must appear in under 50~ms.
        A new request comes with every keypress.
        \textbf{Verdict: real-time.}
        Deploy the model as a REST API (e.g.\ FastAPI). Each keypress
        sends an HTTP request; the API replies in under 50~ms. The model
        stays loaded in memory at all times.

  \item \textbf{Monthly churn report for account managers.}
        Each month, account managers get a list of clients at high risk.
        They review it in a monthly meeting.
        \textbf{Verdict: batch.}
        Run a scoring job on the last day of every month. Export results
        to a spreadsheet. The manager opens it next morning. A real-time
        API would be wasted here.

  \item \textbf{Match the pattern to the need.}
        If you picked real-time for all three, you would waste engineering
        effort maintaining live APIs for two use cases that never need them.
        If you picked batch for the search bar, the product would not work.
        Always ask: \emph{how fresh must the prediction be?}
\end{steps}
\end{worked}

\subsection{Comparing the two patterns side by side}

\medskip
\noindent
\begin{adjustbox}{max width=\linewidth}
\begin{tabular}{@{}p{2.8cm}p{5.2cm}p{5.2cm}@{}}
\toprule
\textbf{Property} & \textbf{Batch} & \textbf{Real-time} \\
\midrule
When predictions run  & Scheduled (hourly, daily, weekly) & On each incoming request \\
Latency               & Minutes to hours; result consumed later & Milliseconds \\
Infrastructure        & Spark, Airflow, a cron job & A persistent API server (FastAPI, TensorFlow Serving) \\
Cost profile          & Lower --- runs when compute is cheap & Higher --- server always on \\
Data freshness        & Stale by design (last batch window) & Fresh at request time \\
Typical use cases     & Emails, reports, segmentation & Fraud, search, real-time recommendations \\
\bottomrule
\end{tabular}
\end{adjustbox}
\medskip

\housefig{figures/fig_batch_vs_realtime.png}{%
  Batch vs.\ real-time serving side by side. Left: a scheduled job stores
  results for later use. Right: a live API answers every request in
  milliseconds. Neither is always better --- pick by latency need.}

\begin{watchout}
\begin{itemize}
  \item \textbf{Batch latency is a hard constraint.}
        If the batch job runs at midnight, predictions are up to 24~hours old.
        If the business needs fresher data, shorten the batch window or
        switch to real-time.
  \item \textbf{Real-time has cold-start costs.}
        A server that is always on costs money even when idle.
        For low-traffic cases, a serverless function (AWS Lambda) can serve
        predictions without a 24/7 server.
        The trade-off: the first request after an idle period is slow
        (that delay is called the \emph{cold start}).
  \item \textbf{Mixing patterns is common.}
        Many systems use batch to pre-compute features and real-time to score
        them. The batch job runs nightly to build feature vectors; the live
        API reads those vectors from a feature store at request time.
        This is called a \emph{lambda architecture}.
  \item \textbf{Real-time is harder to debug.}
        A batch job writes a log file you can inspect. A live API at
        1,000~requests/second produces a flood of logs. Build observability
        (metrics, tracing, alerting) from day one.
\end{itemize}
\end{watchout}

\noindent\textbf{Where this shows up in ML.}
This is one of the most-asked questions in ML system design interviews:
``Would you use batch or real-time serving for \ldots?''
Almost every major ML platform exposes both patterns. \textbf{AWS SageMaker}
has \emph{Batch Transform} (batch) and \emph{Real-Time Inference Endpoints}.
\textbf{GCP Vertex AI} has \emph{Batch Prediction} and \emph{Online Prediction}.
\textbf{TensorFlow Serving} and \textbf{Triton Inference Server} are common
real-time engines. \textbf{Apache Spark} and \textbf{Airflow} are common batch
engines. Knowing which to use, and when to combine them, is a core MLOps skill.

\begin{keytake}
Batch prediction runs on a schedule and stores results for later.
Real-time prediction keeps a live API that answers in milliseconds.
Choose by asking: \emph{how fresh must the prediction be when needed?}
If hours are fine, use batch. If seconds matter, use real-time.
\end{keytake}

% =============================================================================
\section*{Glossary \& symbols}
\addcontentsline{toc}{section}{Glossary \& symbols}
% =============================================================================

\noindent\textbf{Key terms.}
\begin{description}[leftmargin=9.5em,style=nextline,font=\normalfont\bfseries\color{navy}]
  \item[event] a small message saying ``something just happened.''
        Published by a producer; consumed by one or more subscribers.
  \item[event bus] the shared channel that receives and routes events.
        Real tools: Kafka, RabbitMQ, Amazon SQS.
  \item[producer] a service that generates and publishes events.
        Also called a \emph{publisher}.
  \item[consumer] a service that subscribes to and processes events.
        Also called a \emph{subscriber}.
  \item[decoupling] the property that services do not call each other
        directly. Each can change, scale, or fail independently.
  \item[design pattern] a named, reusable solution to a recurring
        engineering problem. A template, not code you copy.
  \item[model registry] a central store for ML model artifacts, versions,
        metadata, and lineage. It bridges experimentation and production.
  \item[model lineage] the full chain from raw data to deployed model:
        which data, code, and hyperparameters produced this artifact.
  \item[alias] a human-readable name (e.g.\ \texttt{production}) that
        points to a specific model version. Promotion means reassigning
        the alias.
  \item[batch prediction] running a model on a large dataset at once,
        on a schedule, storing results for later use.
  \item[online prediction] running a model via a live API that answers
        each request in milliseconds.
  \item[latency] the time between a request arriving and the response
        returning. The key metric for real-time serving.
  \item[cold start] the delay on the first request after an idle period
        (common in serverless real-time serving).
  \item[lambda architecture] a system that combines batch and real-time:
        batch for heavy offline work, real-time for fast online lookups.
\end{description}

\medskip
\noindent\textbf{Tools at a glance.}
\begin{center}\small
\begin{tabular}{@{}p{3.8cm}p{9.0cm}@{}}
\toprule
\textbf{Tool} & \textbf{Role in this session} \\
\midrule
Apache Kafka     & Event bus for high-throughput event-driven systems \\
RabbitMQ         & Message broker; flexible routing and task queues \\
Amazon SQS       & Managed cloud message queue \\
MLflow           & Experiment tracking and model registry (open-source) \\
DVC              & Data and model version control (Git-compatible) \\
Weights \& Biases & Experiment tracking and model registry (cloud) \\
Apache Spark     & Distributed batch processing engine \\
Apache Airflow   & Workflow orchestrator for batch pipelines \\
FastAPI          & Python web framework; common for real-time serving \\
\bottomrule
\end{tabular}
\end{center}

\end{document}
